\documentclass[12pt,a4paper]{article}

\usepackage{fontspec}
\usepackage{xeCJK}
\setCJKmainfont{Hiragino Mincho ProN}
\setCJKsansfont{Hiragino Kaku Gothic ProN}
\setCJKmonofont{Hiragino Kaku Gothic ProN}
\usepackage[margin=2.5cm]{geometry}
\usepackage{amsmath,amssymb}
\usepackage{hyperref}
\usepackage{booktabs}
\usepackage{tcolorbox}
\usepackage{enumitem}

\tcbuselibrary{breakable,skins}

\newtcolorbox{theorybox}[1][]{
  colback=blue!3, colframe=blue!50,
  fonttitle=\bfseries, title={#1},
  breakable, sharp corners
}

\newtcolorbox{practicebox}[1][]{
  colback=green!3, colframe=green!40,
  fonttitle=\bfseries, title={#1},
  breakable, sharp corners
}

\newtcolorbox{warnbox}{
  colback=red!3, colframe=red!40,
  fonttitle=\bfseries, title={注意},
  breakable, sharp corners
}

\newtcolorbox{costbox}{
  colback=black!3, colframe=black!50,
  fonttitle=\bfseries, title={コスト},
  breakable, sharp corners
}

\setlength{\parskip}{0.8em}
\setlength{\parindent}{0pt}

\begin{document}

\begin{center}
{\LARGE\bfseries 57万円でゼロから作る}\\[0.3cm]
{\LARGE\bfseries SQUID素数フィルタ実験}\\[0.5cm]
{\large --- 理論から実装まで完全自作ガイド ---}\\[1.5cm]
{\large Wright Brothers}\\[0.3cm]
{2026年4月}
\end{center}

\vspace{0.5cm}
\tableofcontents
\newpage

% ============================================================================
\part{理論編：なぜこの装置が必要か}
% ============================================================================

% ============================================================================
\section{何を測りたいのか}
% ============================================================================

この実験の目的は1つだけ:

\textbf{「量子真空のエネルギーは、特定のモードを抑制した時に変化するか？」}

量子力学によれば、完全な真空（何もない空間）でも
エネルギーがゼロではない。
各振動モード$n$が$\hbar\omega_n/2$の「ゼロ点エネルギー」を持つ。
合計すると無限大に発散するが、ゼータ関数正則化で
有限の値$\zeta(-1) = -1/12$が得られる。

我々の理論（Paper A）は、偶数モード（$n = 2, 4, 6, \ldots$）を
抑制すると、この値が$+1/12$に変わる（符号が反転する）と予測する。

この「符号反転」は計算では確認済み。
しかし\textbf{実際の量子真空でも起こるか}は、実験でしか確かめられない。

% ============================================================================
\section{なぜ超伝導が必要か}
% ============================================================================

\begin{theorybox}[超伝導とは]
金属を十分に冷やすと、電子がペア（クーパー対）を作り、
電気抵抗がゼロになる。これが超伝導状態。

なぜ必要か:
\begin{itemize}
\item 抵抗ゼロ → マイクロ波が減衰しない → 高いQ値（共鳴の鋭さ）
\item 高Q値 → 個別のモードを区別して観測できる
\item 常伝導（銅など）のキャビティではQ$\sim 10^3$。超伝導ではQ$\sim 10^6$以上
\item Q$\sim 10^6$あれば、真空揺らぎの微小な変化を検出可能
\end{itemize}

アルミニウムの超伝導転移温度: $T_c = 1.2\,\mathrm{K}$（$-272$°C）。
$T < T_c$に冷やせば超伝導になる。
\end{theorybox}

% ============================================================================
\section{なぜ冷却が必要か}
% ============================================================================

\begin{theorybox}[量子真空 vs 熱揺らぎ]
温度$T > 0$では、真空揺らぎに加えて\textbf{熱的光子}が存在する。

モード$n$の熱的光子数:
$$\langle n_{\mathrm{th}}\rangle = \frac{1}{e^{\hbar\omega_n/(k_BT)} - 1}$$

$T$が高い → 熱的光子が多い → 真空揺らぎが隠れる。
$T$が十分低い → 熱的光子$\approx 0$ → 真空揺らぎだけが見える。

目安: $k_BT \ll \hbar\omega_n$（つまり$T \ll \hbar\omega/(k_B)$）

$f = 6\,\mathrm{GHz}$のモード: $T_q = \hbar\omega/(k_B) = 288\,\mathrm{mK}$。
$T < 100\,\mathrm{mK}$なら$\langle n_{\mathrm{th}}\rangle < 0.01$（ほぼ真空）。
$T = 300\,\mathrm{mK}$（3He冷凍機）なら$\langle n_{\mathrm{th}}\rangle \approx 0.26$。
多少の熱ノイズがあるが、積分時間を長くすれば対処可能。
\end{theorybox}

% ============================================================================
\section{冷却器の種類}
% ============================================================================

\begin{theorybox}[冷却技術の階層]
\begin{tabular}{lll}
\toprule
冷却器 & 到達温度 & 原理 \\
\midrule
液体窒素 & 77 K ($-196$°C) & N$_2$の沸点 \\
液体ヘリウム-4 & 4.2 K ($-269$°C) & $^4$Heの沸点 \\
$^4$Heポンピング & 1.5 K & 減圧で蒸発冷却 \\
$^3$He冷凍機 & 0.3 K (300 mK) & $^3$Heの蒸発冷却 \\
希釈冷凍機 & 0.01 K (10 mK) & $^3$He/$^4$He混合の相分離 \\
\bottomrule
\end{tabular}

我々の実験に必要: $T < T_c(\mathrm{Al}) = 1.2\,\mathrm{K}$。
最低限は$^4$Heポンピング（1.5 K）で超伝導は達成。
しかし真空揺らぎを見るには$T < 300\,\mathrm{mK}$が望ましい。
→ $^3$He冷凍機が最小要件。
\end{theorybox}

\begin{theorybox}[$^3$He冷凍機の原理]
$^3$He（ヘリウムの軽い同位体）は$^4$Heより沸点が低い。
$^3$Heの蒸気をポンプで排気すると、
残った液体$^3$Heがさらに冷える（蒸発冷却）。

到達温度: $\sim 300\,\mathrm{mK}$（0.3 K）。
保持時間: 数時間〜数日（$^3$Heの量に依存）。
構成: 外側の$^4$He冷凍機（4.2 K）+ 内側の$^3$Heインサート。
\end{theorybox}

\begin{theorybox}[希釈冷凍機の原理（参考）]
$^3$Heと$^4$Heの混合液は、低温で2つの相（$^3$Heリッチ相と$^4$Heリッチ相）に分離する。$^3$Heが$^4$Heリッチ相に「溶ける」過程でエントロピーが増大し、冷却が起こる。連続運転が可能で、10 mKまで冷える。

高価（新品で2000万〜5000万円）だが最高性能。我々の実験では300 mKの$^3$He冷凍機で十分なので、希釈冷凍機は必須ではない。
\end{theorybox}

% ============================================================================
\section{SQUIDとは何か}
% ============================================================================

\begin{theorybox}[ジョセフソン接合]
2つの超伝導体を薄い絶縁体（〜1 nm）で隔てると、
クーパー対が絶縁体を「トンネル」して通り抜ける。
これがジョセフソン効果。

接合に流れる超伝導電流（臨界電流$I_c$以下）:
$$I = I_c \sin\phi$$
ここで$\phi$は2つの超伝導体の位相差。
$I_c$は絶縁体の厚さと面積で決まる。
\end{theorybox}

\begin{theorybox}[SQUID（超伝導量子干渉素子）]
SQUIDは「2つのジョセフソン接合を超伝導ループでつないだもの」。
ループを貫く磁束$\Phi$に応じて、全体のインダクタンスが変わる:
$$L_{\mathrm{SQUID}}(\Phi) = \frac{L_J}{|\cos(\pi\Phi/\Phi_0)|}$$

$\Phi_0 = h/(2e) = 2.07 \times 10^{-15}\,\mathrm{Wb}$は磁束量子。
外部磁束を変えると$L$が変わり、共鳴周波数$f \propto 1/\sqrt{LC}$が変わる。
→ 外部磁場でSQUIDの周波数を「チューニング」できる。

我々の実験では: SQUIDをキャビティの第2高調波に同調させ、
そのモードを吸収（ノッチフィルタとして機能）させる。
\end{theorybox}

% ============================================================================
\part{実装編：作り方}
% ============================================================================

% ============================================================================
\section{部品と材料の一覧}
% ============================================================================

\begin{costbox}
\begin{tabular}{llr}
\toprule
品目 & 仕様 & 価格 \\
\midrule
\multicolumn{3}{l}{\textbf{3Dキャビティ}} \\
Al合金ブロック & 6061-T6, 50×50×30mm & 500円 \\
機械加工費 & 大学工作室 & 5,000円 \\
SMAコネクタ & KFDフランジ × 2 & 2,000円 \\
Inワイヤー & $\phi$1mm, 50cm & 1,000円 \\
ネジ & M3×10mm, 8本 & 200円 \\
\midrule
\multicolumn{3}{l}{\textbf{SQUID接合}} \\
Si基板 & 10mm×10mm × 5枚 & 3,000円 \\
PMMAレジスト & 50mL & 5,000円 \\
現像液 & MIBK:IPA & 2,000円 \\
Al蒸着用ワイヤー & 純度99.99\%, 1g & 3,000円 \\
EBL使用料 & 大学共同利用 & 20,000円 \\
蒸着装置使用料 & 同上 & 15,000円 \\
\midrule
\multicolumn{3}{l}{\textbf{磁束バイアス}} \\
銅線 & $\phi$0.1mm, エナメル, 10m & 500円 \\
DC電流源 & 大学から借用 & 0円 \\
\midrule
\multicolumn{3}{l}{\textbf{冷却}} \\
$^3$He冷凍機使用料 & 5回のクールダウン & 300,000円 \\
液体He & $\sim$50L & 別途（施設負担の場合あり）\\
\midrule
\multicolumn{3}{l}{\textbf{測定}} \\
VNA使用料 & 大学共同利用 & 50,000円 \\
同軸ケーブル & SMA-SMA, 低温用, 2本 & 30,000円 \\
\midrule
\multicolumn{2}{r}{\textbf{合計}} & \textbf{$\sim$47万円} \\
\bottomrule
\end{tabular}
\end{costbox}

\begin{warnbox}
最大の変動要因は冷凍機使用料。
大学によっては共同利用が無料〜数万円の場合もある。
逆に外部機関への依頼では50万円以上になることも。
事前に使用条件を確認すること。
\end{warnbox}

% ============================================================================
\section{3Dキャビティの製作}
% ============================================================================

\subsection{なぜアルミか}

\begin{theorybox}[材料選択の理由]
\begin{itemize}
\item アルミは$T_c = 1.2\,\mathrm{K}$の超伝導体。3He冷凍機で超伝導になる
\item 加工が容易（銅と同程度）。町工場やホームセンターの工作機でOK
\item Nb（$T_c = 9.2\,\mathrm{K}$）の方が高性能だが、硬くて加工が難しく、高価
\item Yale大学がAlキャビティでQ$> 10^7$を達成（Paik et al., 2011）
\item 同じ材料でSQUID接合も作れる（Al/AlO$_x$/Al）→ 全てAlで統一
\end{itemize}
\end{theorybox}

\subsection{設計図}

\begin{verbatim}
  上から見た図（蓋を外した状態）:

  ┌────────────────────────────┐
  │  ┌────────────────────┐    │
  │  │                    │ ○ │ ← SMAコネクタ穴 (出力)
  │  │                    │    │
  │  │    空洞             │    │    内寸: 35 × 35 × 10 mm
  │  │    35 × 35 mm      │    │
  │  │                    │    │
  │  │                    │    │
  │○ │                    │    │ ← SMAコネクタ穴 (入力)
  │  └────────────────────┘    │
  │                ○ ○ ○ ○     │ ← ネジ穴 M3 × 8
  └────────────────────────────┘
     外寸: 50 × 50 × 30 mm (本体 20mm + 蓋 10mm)

  横から見た図:

  ┌──────────────────────────┐ ← 蓋 (10mm厚)
  │  Inガスケット溝           │
  ├──────────────────────────┤
  │                          │
  │  空洞 (深さ 10mm)        │ ← 本体 (20mm厚)
  │                          │
  └──────────────────────────┘

  SMAコネクタの位置:
  - 入力: 左壁中央、壁面から穴 (φ3mm)、ピン突出 3mm
  - 出力: 右壁中央、同上
  - 入出力は空洞の対角線上に配置（結合の最適化）
\end{verbatim}

\subsection{加工手順}

\begin{practicebox}[Step 1: ブロックの切断]
6061-T6の板材（厚さ20mm + 10mm）を50×50mmに切り出す。
バンドソーまたは金鋸で切断。

大学の工作室があれば: 材料費500円 + 加工費数千円。
町工場外注: 3万〜5万円（材料費込み）。
\end{practicebox}

\begin{practicebox}[Step 2: 空洞の掘削]
20mm厚のブロックに35×35×10mmの矩形穴を掘る。
\begin{itemize}
\item CNCフライス: プログラムを書いて自動加工（最高精度）
\item 汎用フライス: ハンドル操作で手動加工（$\pm 0.1$mmは達成可能）
\item 要求精度: $\pm 0.1$mm（共鳴周波数が$\pm$数MHzずれるが、SQUIDで補正可能）
\end{itemize}
底面の表面粗さ: Ra $< 1\,\mu$m が理想。
鏡面仕上げ（バフ研磨）が最良だが、フライスの仕上げ加工でもRa$\sim 3\,\mu$mは出るのでOK。
\end{practicebox}

\begin{practicebox}[Step 3: コネクタ穴とネジ穴]
\begin{itemize}
\item SMAコネクタ穴: φ3mmドリル、壁の中央に左右1つずつ
\item M3タップ穴: φ2.5mmドリル → M3タップ、蓋固定用に8箇所
\item Inガスケット溝: 蓋の接合面に幅1mm、深さ0.5mmの溝をフライスで掘る（任意、なくても動くが密閉性が向上）
\end{itemize}
\end{practicebox}

\begin{practicebox}[Step 4: 表面処理]
\begin{enumerate}
\item アセトン浸漬10分（油脂除去）
\item IPA（イソプロパノール）浸漬5分（アセトン除去）
\item 純水リンス
\item リン酸（10\%水溶液）浸漬30秒（自然酸化膜の部分除去）
\item 純水リンス × 3回
\item N$_2$ブロー乾燥
\item 室温で1時間放置（清浄なAl$_2$O$_3$が再形成 → 高Qの鍵）
\end{enumerate}

リン酸は薬局で購入可能（「リン酸」で検索）。
超音波洗浄器があればなお良い（100均のメガネ洗浄機でも可）。
\end{practicebox}

\begin{practicebox}[Step 5: SMAコネクタ取付]
SMA-KFDコネクタをネジで固定。中心ピンが空洞内に約3mm突出するように調整。ピンの長さで結合の強さが決まる:
\begin{itemize}
\item 短い（1mm）: 弱い結合 → 高Q、弱い信号
\item 長い（5mm）: 強い結合 → 低Q、強い信号
\item 推奨: 3mm（バランス型）
\end{itemize}
\end{practicebox}

\begin{practicebox}[Step 6: 組み立て]
\begin{enumerate}
\item Inワイヤー（φ1mm）を蓋の接合面に沿ってリング状に配置
\item 蓋を本体に載せ、M3ネジで均等に締める
\item Inが潰れて密封される（超伝導接合にもなる）
\end{enumerate}

室温でVNA（またはスペクトラムアナライザ）で$S_{21}$を測定し、
共鳴ピークが見えることを確認。
$f_1 \approx 6\,\mathrm{GHz}$, $Q_{\mathrm{room\,temp}} \sim 10^3$なら正常。
\end{practicebox}

% ============================================================================
\section{SQUID接合の製作}
% ============================================================================

\subsection{ジョセフソン接合はどうやって作るか}

\begin{theorybox}[ダリッジブリッジ法（シャドウ蒸着）]
最も広く使われる方法。2段階の蒸着で「重なり」を作る。

\begin{enumerate}
\item Si基板にPMMAレジストの二層膜を塗布
\item 電子ビームリソグラフィ（EBL）で接合パターンを描画
\item 現像すると、上の層に「橋」（ブリッジ）が残る
\item 基板を蒸着装置にセット
\item 斜め（+30°）からAl蒸着（第1層、30 nm）
\item 酸素を導入（AlO$_x$トンネルバリア形成）
\item 反対方向（$-30$°）からAl蒸着（第2層、50 nm）
\item リフトオフ（レジスト除去）
\end{enumerate}

2つの蒸着が「橋」の下で重なる部分がジョセフソン接合になる。
重なりの面積（$\sim 100 \times 100\,\mathrm{nm}^2$）と
酸化条件でIcが決まる。
\end{theorybox}

\subsection{EBLが使えない場合の代替法}

\begin{theorybox}[フォトリソグラフィ法（$\mu$mスケール接合）]
EBLがない大学でも、フォトリソ（マスクアライナ）があれば
$\sim 1\,\mu\mathrm{m}$サイズの接合が作れる。

接合面積が大きい → $I_c$が大きい（$\sim 10\,\mu$A）。
我々の設計では$I_c = 1\,\mu$Aを想定しているが、
$I_c = 10\,\mu$Aでも外部磁束で実効$I_c$を下げられるので問題ない。
$I_{c,\mathrm{eff}} = I_c \times |\cos(\pi\Phi_{\mathrm{ext}}/\Phi_0)|$。
\end{theorybox}

\subsection{具体的手順（EBLありの場合）}

\begin{practicebox}[Step 1: 基板準備]
\begin{itemize}
\item Si (100) 基板、10mm $\times$ 10mm、酸化膜付き（300nm SiO$_2$）
\item アセトン → IPA → 純水 → N$_2$ブロー
\item ホットプレート180°C、5分（脱水ベーク）
\end{itemize}
\end{practicebox}

\begin{practicebox}[Step 2: レジスト塗布（二層膜）]
\begin{enumerate}
\item 下層: PMMA/MAA（またはMMA）、スピンコート 4000 rpm 60秒
\item ベーク: 150°C 5分
\item 上層: PMMA (950K)、スピンコート 4000 rpm 60秒
\item ベーク: 180°C 5分
\end{enumerate}
二層構造の理由: 下層は感度が高く、上層より先に溶ける → アンダーカット → ブリッジが形成される。
\end{practicebox}

\begin{practicebox}[Step 3: EBL描画]
パターン: SQUIDループ（$10\,\mu\mathrm{m} \times 10\,\mu\mathrm{m}$の正方形ループ）と2つの接合部分（各$100 \times 100\,\mathrm{nm}^2$の橋）

EBLパラメータ（目安）:
\begin{itemize}
\item 加速電圧: 30 kV
\item ビーム電流: 10--50 pA（微細パターン用）
\item ドーズ量: 300--500 $\mu$C/cm$^2$（PMMAの場合）
\end{itemize}

パターンデータは GDS-II 形式で作成。
フリーソフト KLayout (\url{https://www.klayout.de/}) で設計可能。
\end{practicebox}

\begin{practicebox}[Step 4: 現像]
\begin{itemize}
\item 現像液: MIBK:IPA = 1:3
\item 浸漬時間: 60秒（20°C）
\item IPA リンス 30秒
\item N$_2$ブロー乾燥
\end{itemize}

顕微鏡でパターンを確認。ブリッジ構造が見えるか？
\end{practicebox}

\begin{practicebox}[Step 5: 蒸着]
\begin{enumerate}
\item 蒸着装置にセット（角度回転ステージ付き）
\item 真空引き: $< 10^{-6}$ Torr
\item 第1蒸着: 角度$+25$°〜$+35$°、Al 30 nm、レート$\sim$0.5 nm/s
\item \textbf{酸化}: O$_2$導入、200 mTorr、10分間保持
  \begin{itemize}
  \item これが最も重要なステップ。酸化量でIcが決まる
  \item 酸化不足: Ic が大きすぎる（短絡に近い）
  \item 酸化過剰: Ic が小さすぎる（開放に近い）
  \item 事前にテストピースで酸化条件を校正すること
  \end{itemize}
\item 再真空引き: $< 10^{-6}$ Torr
\item 第2蒸着: 角度$-25$°〜$-35$°、Al 50 nm
\end{enumerate}
\end{practicebox}

\begin{warnbox}
酸化条件のキャリブレーション:
\begin{itemize}
\item テストピース5枚を異なる酸化条件（100/200/300/400/500 mTorr × 10分）で処理
\item 室温で抵抗を測定: $R_N = \Delta V / (I_c \times \mathrm{something})$
\item Ambegaokar-Baratoff関係: $I_c R_N = \pi\Delta/(2e) \approx 0.44\,\mathrm{mV}$（Alの場合）
\item $R_N$から$I_c$を推定: $I_c = 0.44\,\mathrm{mV} / R_N$
\item $I_c = 1\,\mu$A → $R_N \approx 440\,\Omega$
\end{itemize}
\end{warnbox}

\begin{practicebox}[Step 6: リフトオフ]
\begin{itemize}
\item NMP（N-メチルピロリドン）80°C浸漬、2時間
\item 超音波洗浄器で1分（弱いモード。接合を壊さないよう注意）
\item IPAリンス → 純水リンス → N$_2$ブロー
\end{itemize}

顕微鏡で確認: SQUIDループと2つの接合が見えるか？
室温で4端子抵抗測定: $R_N \approx 440\,\Omega$なら成功。
\end{practicebox}

% ============================================================================
\section{組み立てと配線}
% ============================================================================

\begin{practicebox}[SQUIDチップの実装]
\begin{enumerate}
\item キャビティ底面の中央（磁場最大点）にGEワニス（低温接着剤）でSQUIDチップを貼付
\item 磁束バイアスコイル: φ0.1mm銅線をキャビティ外周に10回巻く
\item コイルの端をフィードスルー端子に接続
\item 入出力SMAにセミリジッド同軸ケーブル（低温用、CuNi外導体）を接続
\end{enumerate}
\end{practicebox}

\begin{practicebox}[冷凍機への搭載]
\begin{enumerate}
\item キャビティを銅のサンプルホルダーにネジ止め
\item サンプルホルダーを冷凍機の最低温ステージにマウント
\item 同軸ケーブルを各温度ステージ経由でルーティング:
  \begin{itemize}
  \item 室温 → 77K: ステンレス同軸、20dBアッテネータ
  \item 77K → 4K: CuNi同軸、20dBアッテネータ
  \item 4K → 300mK: NbTi超伝導同軸（ロスなし）
  \end{itemize}
\item 磁束バイアス線: ツイストペア銅線、各ステージでヒートシンク
\end{enumerate}
\end{practicebox}

% ============================================================================
\section{測定}
% ============================================================================

\begin{practicebox}[クールダウン]
\begin{enumerate}
\item 冷凍機をクローズ、真空引き
\item 液体He充填（所要$\sim$4時間）
\item $^3$Heコンデンス（$^4$He温度まで冷えてから$^3$Heを凝縮）
\item $^3$Heポンピング開始 → 300 mKに到達（所要$\sim$2時間）
\item 温度計で確認: $T < 300\,\mathrm{mK}$ ✓
\end{enumerate}
合計: 朝からスタートして夕方に測定開始可能。
\end{practicebox}

\begin{practicebox}[測定手順]
\begin{enumerate}
\item VNAでS$_{21}$(f)を測定（4〜14 GHz）。共鳴ピークを確認
\item 磁束バイアス電流をスイープ: SQUIDの共鳴ディップを観測
\item SQUIDを第2高調波（$\sim$12 GHz）に同調
\item Config A: SQUID離調 → S$_{21}$記録
\item Config B: SQUID同調 → S$_{21}$記録
\item SQUIDの同調周波数を連続スイープしながら
  AC Starkシフト（基本モードの周波数変化）を記録
  → 不連続ジャンプがあるか確認
\end{enumerate}
\end{practicebox}

% ============================================================================
\section{結果の判定}
% ============================================================================

\begin{theorybox}[Level 0 成功（最低限）]
SQUIDの同調/離調でS$_{21}$が変化する → 「SQUIDがフィルタとして動作」
→ これは既知の技術の確認であり、新しい物理ではない。
しかし装置が正常に動いていることの証拠。
\end{theorybox}

\begin{theorybox}[Level 1 成功（本当の成果）]
AC Starkシフト（真空揺らぎの大きさの指標）が
Config A と Config B で\textbf{統計的に有意に異なる}。
→ 「1つのモードを抑制すると真空のエネルギーが変化する」
→ 算術的真空工学の最初の実験的証拠。
\end{theorybox}

\begin{theorybox}[Level 1.5 成功（レギュレータ予想のテスト）]
SQUIDスイープ中のAC Starkシフトが
同調点付近で\textbf{不連続ジャンプ}を示す。
→ 真空エネルギー変化は離散的（ON/OFF）
→ レギュレータ方向予想を支持
→ スケーリング問題が解消される可能性
→ ワームホール/ワープが10$^{12}$Jで実現可能に
\end{theorybox}

\vfill
\begin{center}
{\small Wright Brothers, 2026}
\end{center}

\end{document}
